\chapter{The Runtime Decision Model}\label{ch09:decision}

\section{The last layer}\label{ch09:intro}

Every mechanism developed so far produces or protects \emph{structure}: references that cannot conflate namespaces, boundaries that cannot contradict themselves, encodings that cannot pun types, compilations that cannot dodge replay, registries that cannot shed their bytes. The runtime decision model is where all of that structure arrives at a question none of it can answer alone: \emph{may this action happen?}

The kernel's answer machinery has a deliberately strange shape. It distinguishes two kinds of question and gives them categorically different treatments. For \emph{nonclaim research evidence}---a request to accept a legacy compatibility evaluation as structurally sound, local, and effect-free---the model performs the most demanding comparison it can construct: a full fresh replay of the underlying compilation, with every exact field and digest of the presented evidence validated against the replay, and a structural pass granted only on total agreement. For \emph{claim-bearing intent}---any action that would assert, publish, license, or otherwise promote a claim---the model performs no comparison at all, because no comparison could help: the answer is a constant, and the constant is no.

This chapter develops the machinery in three movements. First, the typed action vocabulary: why every action is an explicit enum triple rather than a string, and what the descriptor that carries it binds. Second, the structural replay comparison: what the gate actually checks before a nonclaim pass, denial by denial. Third, the two constant-false invariants and the design constraints any future authorization implementation must satisfy---including the argument, made in full, for why ``there is no code path to yes'' is a strictly stronger safety property than ``the policy says no.''

As everywhere in this monograph, the object under discussion is a Ring-4 research candidate. The decision model's behaviors are software behaviors, test-evidenced at the Gate~4 receipt under exact fixtures; nothing in this chapter is, or licenses, an authority claim---indeed the chapter's central point is that the model itself refuses to let its outputs be read as one.

\section{Typed actions: the end of stringly-typed intent}\label{ch09:actions}

\subsection{The three enums}\label{ch09:enums}

Every action presented to the decision model is described by three explicit enumerations, with no string-typed alternative path.

\begin{definition}[Action vocabulary]\label{ch09:def:vocab}
An action's \emph{operation} is one of \texttt{LEGACY\_COMPATIBILITY\_EVALUATION}, \texttt{MODEL\_INFERENCE}, or \texttt{PUBLISH\_CLAIM}. Its \emph{target} is one of \texttt{LOCAL\_RING4\_RESEARCH\_ARTIFACT}, \texttt{EXTERNAL\_SYSTEM}, or \texttt{PUBLIC\_CLAIM\_SURFACE}. Its \emph{effect class} is one of \texttt{PURE\_LOCAL\_COMPUTE}, \texttt{CLAIM\_BEARING}, or \texttt{EXTERNAL\_SIDE\_EFFECT}. The three axes are independent, and the decision model receives all three explicitly.
\end{definition}

Notice first what the vocabulary encodes: the enumerations name not only what the kernel supports but what it \emph{refuses}. Only one operation (\texttt{LEGACY\_COMPATIBILITY\_EVALUATION}), one target (\texttt{LOCAL\_RING4\_RESEARCH\_ARTIFACT}), and one effect class (\texttt{PURE\_LOCAL\_COMPUTE}) can ever appear in a passing decision in v0.1; the other six values exist so that the refusals are typed. A model-inference action, an external-system target, a claim-bearing effect---each is representable precisely so that presenting one produces a specific, enumerable denial rather than an unparseable surprise. The vocabulary is a map of the kernel's frontier: everything on it is either the single admitted interior point or a named point on the boundary.

Why enums rather than strings? The question sounds stylistic; it is not. Consider what a stringly-typed action interface---\texttt{decide("legacy\_compat\_eval", ...)}---costs an evidence kernel, failure by failure.

\emph{Silent vocabulary growth.} With strings, the set of possible actions is whatever callers happen to spell, and it grows by typo, by convention drift, and by helpful extension, with no review event marking the growth. The decision logic inevitably develops a default arm---\texttt{else: deny}, if one is lucky; \texttt{else: log and continue}, in the historical worst case---and new action kinds are born into the default arm, their handling determined by accident. With enums, adding an action kind is an API change: it appears in a diff, it breaks exhaustiveness expectations, it demands a decision-model arm, and it cannot be introduced by a caller at all. The set of things the kernel can be asked is closed under review, which for a kernel whose whole purpose is refusing unauthorized promotion is not a convenience but a precondition.

\emph{Aliasing and normalization ambiguity.} Strings invite the question of which strings are the same action: case variants, hyphen and underscore variants, trailing whitespace, Unicode confusables. Every normalization rule adopted to answer it becomes attack surface---the encoding chapter's identifier-profile discipline exists precisely because such rules have been exploited elsewhere---and every rule not adopted becomes a source of two-spellings-one-intent bugs. Enum members have no spellings at the decision layer; equality is identity of the member, and the serialized form is a single canonical string controlled by the kernel, not the caller.

\emph{Unauditable comparisons.} A string comparison in decision logic tells the reviewer nothing about the universe of values it discriminates: \texttt{if op == "publish\_claim"} is one arm of an unknowably large conditional. Typed comparisons make the universe explicit, so the reviewer of the decision function can verify---locally, without global search---that every member of every axis has a determined fate. The v0.1 decision logic has exactly this shape: each enum axis is matched member-by-member, with the unsupported members producing their named denial codes, and the exhaustiveness is checkable by inspection.

The kernel hardens the typing one step further than most codebases: the decision path checks that each supplied value's runtime type is \emph{exactly} the enum type---not a subclass, not a value-equal string. The enums inherit from string for serialization convenience, which means a bare string equal to a member's value would compare equal under ordinary equality; the exact-type check ensures that a caller who bypasses the enum and passes the raw string is rejected rather than accidentally honored. This is the semantic-identity chapter's exact-type rule, applied to intent: equality-compatible impostors are not admitted as identity-bearing values, whether they impersonate a reference, a boundary, or an action kind.

\subsection{The action descriptor}\label{ch09:descriptor}

The three enum values travel inside a frozen, slotted descriptor that binds them to the identities of everything the action concerns. A descriptor carries: its own version string (rejected unless it is the single known version); the operation, target, and effect-class members (exact-typed as above); the SHA-256 identities of the source under evaluation, the compilation result, the model (if any), and the environment; optional claim identifiers; the descriptor's own canonical digest; and the research boundary, which must be exactly Ring 4 / \textsc{open} / \textsc{not licensed}.

Construction is where the descriptor earns its keep, because every coherence rule is enforced there, and an incoherent descriptor therefore never exists as a value. The digest fields must each be syntactically valid SHA-256 identifiers. The claim fields obey a pairing rule: \texttt{claim\_id} and \texttt{claim\_version} are supplied together or not at all, since a claim identifier without a version is exactly the ambiguity Section~\ref{ch09:future} shows authorization cannot tolerate, and a version without an identifier is noise. The claim identifier, when present, must satisfy the ASCII identifier profile---review-visible machine identifiers never carry unrestricted Unicode. And two cross-field rules tie intent to metadata: a \texttt{PURE\_LOCAL\_COMPUTE} action \emph{cannot} carry claim metadata (a ``pure local'' action naming a claim is either confused or laundering, and both are rejected), while a \texttt{CLAIM\_BEARING} action \emph{must} carry exact claim metadata (a claim-bearing action that declines to say which claim is unreviewable by construction). Subclassing the descriptor is rejected outright, closing the equality-override impersonation channel.

Two further details complete the picture. First, the model-identity field: the legacy compatibility path contains no model invocation, and rather than leaving the field empty---emptiness being exactly the kind of unclaimed territory that later fills in silently---the kernel defines a canonical \emph{no-model} identity, the digest of a tagged structure stating that the deterministic frozen compiler path invokes no model. An action on the legacy path must carry precisely that digest; any other value denies with \texttt{DENY\_UNBOUND\_OR\_UNEXPECTED\_MODEL\_IDENTITY}. Absence, made exact, becomes checkable. Second, the self-digest: the descriptor's \texttt{action\_sha256} is the canonical digest of its own payload, verified at construction and re-verified at assessment (\texttt{DENY\_CORRUPT\_ACTION\_DESCRIPTOR} on mismatch). As with every self-digest in this kernel, its scope is corruption detection: it proves the descriptor is the descriptor it was when built, and deliberately proves nothing about who built it.

\begin{ruleenv}[No stringly-typed actions]\label{ch09:rule:typed-actions}
Action intent enters the decision model only as exact-typed enum members bound inside a version-checked, digest-verified, boundary-carrying descriptor whose coherence rules are enforced at construction. There is no string-accepting decision path, no default arm that admits unrecognized intent, and no representable descriptor whose intent and claim metadata contradict each other.
\end{ruleenv}

\section{The structural replay comparison}\label{ch09:replay}

\subsection{What a nonclaim pass requires}\label{ch09:pass-requirements}

Nonclaim research evidence earns, at most, a \emph{structural pass}: a decision whose codes state that a Ring-4, pure-local-compute action is structurally admissible, accompanied---always---by the two limitation codes recording the non-sandbox boundary and by the standing code that authorization is not granted. The path to that pass is the most heavily checked route in the kernel, and it is worth walking in full, because its design principle---\emph{every exact compatibility field and digest must validate, against a fresh replay, before anything passes}---is applied without exception.

The evidence presented to the gate is a compilation assessment: a typed object binding a compiler-validation receipt, an action descriptor, the underlying compilation record, the canonical digest of that compilation's evidence, and its own self-digest. The gate's evaluation proceeds in layers, and each layer's failure produces a distinct denial code, so the decision record identifies exactly where the evidence fell short.

First, \emph{shape and type}: every component must be exactly its expected type---the receipt a real receipt, the action a real descriptor, the compilation a real compilation record, the boundary exactly Ring-4 open---with the component-wise validators re-run (\texttt{DENY\_ASSESSMENT\_TYPE\_MISMATCH} on any failure). Second, \emph{integrity}: the assessment's self-digest must verify (\texttt{DENY\_CORRUPT\_COMPILATION\_ASSESSMENT}), and if the caller supplied an action alongside the evidence, it must match the assessed action field-for-field under canonical serialization (\texttt{DENY\_ACTION\_ASSESSMENT\_MISMATCH})---the gate will not let evidence assessed for one action be spent on another.

Third, and centrally, \emph{replay}: the gate re-executes the entire deterministic-closure protocol of Chapter~\ref{ch07:worker-protocol} on the compilation's retained source and compiler bytes. A replay that fails for any reason denies with \texttt{DENY\_COMPILATION\_REPLAY\_FAILED}; the gate never proceeds on stored evidence alone (Invariant~\ref{ch07:inv:replay}). From the replay it derives a fresh receipt and a fresh evidence digest, and now the comparisons begin: the presented receipt must match the replay-derived receipt in every field under canonical byte comparison (\texttt{DENY\_RECEIPT\_REPLAY\_MISMATCH}); the presented evidence digest must equal the fresh one (\texttt{DENY\_COMPILATION\_EVIDENCE\_MISMATCH}); and the presented assessment's verdict and code tuple must equal what the gate itself computes from the fresh receipt and the action (\texttt{DENY\_ASSESSMENT\_STATE\_MISMATCH})---so even an evidence object whose components are individually honest is denied if its stated conclusion differs from the freshly recomputed one.

Fourth, the \emph{receipt-versus-action findings}, computed on the fresh receipt: the receipt's self-digest; its result validity; both compiler identities against the frozen constants; the source authority; the semantic, syntax, and engine profiles against their frozen strings; the isolation booleans (isolated, environment-ignored, fresh, uninherited---else \texttt{DENY\_NONISOLATED\_PYTHON\_PROCESS}); the false capability fields (a receipt asserting sandboxing denies with \texttt{DENY\_FALSE\_OS\_ISOLATION\_CLAIM}); the \texttt{legacy\_only} flag; the self-hash scope string; and then the action's bindings against the receipt's---source digest (\texttt{DENY\_ACTION\_SOURCE\_MISMATCH}), result digest (\texttt{DENY\_ACTION\_RESULT\_MISMATCH}), environment digest (\texttt{DENY\_ACTION\_ENVIRONMENT\_MISMATCH}), model identity, operation, target, and effect class. Only if every finding is clean does the decision carry the structural-pass code, and the pass is computed, not copied: the gate's verdict is always the verdict of its own fresh computation.

Finally, if a claim registry is supplied alongside the evidence, the gate verifies the registry's raw byte bindings---the exact-byte-backed loading discipline of the registry chapter---before recording the registry's raw digest in the decision (\texttt{DENY\_CLAIM\_REGISTRY\_RAW\_BINDING} on failure). Note carefully what this check is \emph{for}: it validates the registry's integrity as presented context; it does not, and cannot, convert any registry content into permission. That refusal is the subject of the next section.

\begin{invariant}[Total validation before structural pass]\label{ch09:inv:total}
A structural pass is issued only when every exact compatibility field and digest of the presented evidence---receipt fields, action bindings, evidence digest, verdict, and codes---validates against a fresh execution of the full closure protocol. There is no partial pass, no pass on stored evidence, and no pass whose content differs from the gate's own recomputation.
\end{invariant}

\subsection{The assessment as an artifact}\label{ch09:assessment-artifact}

Between the compilation and the gate stands an intermediate artifact worth examining in its own right: the compilation assessment. Its construction path is the model's producer side, and its anatomy explains several gate behaviors that otherwise look like duplication.

An assessment is produced by one admitted function, which takes an exact compilation and an exact action descriptor and---before assessing anything---freshly replays the compilation through the full closure protocol. The receipt it embeds is derived from the \emph{replay}, not from the caller's record; the evidence digest it embeds is the canonical digest of the replay's complete serialization; the verdict and code tuple it embeds are computed from that fresh receipt against the action; and the whole payload is sealed under the assessment's own canonical self-digest and stamped with the replay's Ring-4 boundary. The constructor then re-verifies everything it was just handed: exact types for receipt, action, and compilation; nonempty typed code tuple; valid digests; boundary exactly Ring-4 open; self-digest verification as a construction requirement. An assessment whose parts do not cohere is, like every other incoherent object in this kernel, not a value that exists.

Why then does the gate---which receives this carefully constructed object---replay \emph{again} and recompute the verdict the assessment already contains? Because the gate's trust model does not include the assessment's history. From the gate's standpoint, an assessment is bytes-shaped input from an unauthenticated source: perhaps it came from the admitted constructor a moment ago; perhaps it was deserialized from last month; perhaps it was assembled by something imitating the constructor. The self-digest distinguishes corruption, not provenance---the kernel's constant refrain---so the gate treats the assessment's embedded verdict as a \emph{claim to be checked}, replays the compilation itself, recomputes the expected verdict and codes, and denies on any mismatch (\texttt{DENY\_ASSESSMENT\_STATE\_MISMATCH}) even when every embedded component is internally valid. The assessment is thus best understood not as a credential but as a \emph{falsifiable prediction bundle}: it packages, in exact canonical form, everything the gate will need to re-derive, so that the gate's agreement or disagreement is total and attributable. The apparent duplication is the design: producer-side replay keeps dishonest assessments from being constructed by admitted code; gate-side replay keeps honestly constructed assessments from being trusted merely for having been constructed.

This double-replay structure also fixes the model's freshness semantics without any appeal to clocks. An assessment carries no timestamp that the gate honors; its validity is not a time window but a \emph{world condition}---the bound files still have their bytes, the closure still reproduces---checked by re-execution at the moment of decision. A stale assessment does not expire; it starts failing its replay, with a code that says exactly what drifted. Time-based validity would have been both weaker (a fresh timestamp proves nothing about the world) and more manipulable (clocks are data); the kernel's version of ``recent enough'' is ``still reproducible now,'' which is the only sense of recency a determinism claim can honestly use.

\subsection{The evidence hierarchy}\label{ch09:hierarchy}

The gate accepts more shapes of evidence than the full assessment, and its treatment of the lesser shapes is itself part of the design, because each lesser shape corresponds to a shortcut someone will eventually try to take.

Presenting a \emph{bare compiler-validation receipt}---no assessment, no bound compilation, no replay possibility---does not produce a smaller pass; it produces a denial whose leading code names the mistake: \texttt{DENY\_STRUCTURAL\_RECEIPT\_NOT\_AUTHORIZATION}. The receipt's self-digest is still checked, and a corrupt receipt earns the additional \texttt{DENY\_CORRUPT\_COMPILER\_VALIDATION}, so the denial record remains diagnostically honest---but no field-level validity can compensate for the absence of the replayable compilation. The design point is anti-inflationary: a receipt is the \emph{summary} of evidence, and summaries do not pass gates. If bare receipts earned even a reduced pass, the incentive gradient for every integration would run toward passing summaries around and away from retaining the replayable records, and the kernel's most expensive guarantee---Invariant~\ref{ch07:inv:replay}---would be priced out of practice by its own convenience layer.

Presenting \emph{nothing} denies with \texttt{DENY\_MISSING\_COMPILATION\_ASSESSMENT}; presenting an object of any unexpected type denies with \texttt{DENY\_UNKNOWN\_STRUCTURAL\_EVIDENCE}. The latter is a wider net than it looks: it is the arm that catches every clever wrapper, subclass, duck-typed imitation, and serialized-then-tampered reconstruction that fails the exact-type discipline. The gate does not attempt to salvage value from unrecognized evidence---salvage logic is exactly where lenient parsers grow---and the unknown-evidence denial is deliberately uninformative about \emph{how} the object fell short, because a detailed rejection oracle is a debugging aid for forgers first and reviewers second.

Alongside whatever evidence is presented, a malformed \emph{action} argument (wrong type, subclass, or non-descriptor) adds \texttt{DENY\_UNKNOWN\_ACTION\_DESCRIPTOR}, and a malformed \emph{registry} argument adds \texttt{DENY\_UNKNOWN\_CLAIM\_REGISTRY}; each also forces the structural verdict false regardless of the evidence's own quality. The rule generalizing these arms: any component of the call that is not exactly what it claims to be poisons the whole decision. The gate never quarantines a bad argument and proceeds with the rest, because a decision computed over a partially trusted call is a decision whose provenance no one can state.

\begin{remark}[Denials are cumulative, ordered, and deduplicated]\label{ch09:rem:codes}
Decision codes accumulate: a single call can earn an evidence-tier denial, an action-type denial, and a registry denial together, and the standing not-granted code is appended to every decision without exception. The tuple preserves first-occurrence order and removes duplicates, so identical failure situations produce identical code tuples---which matters because the code tuple participates in the assessment's canonical serialization and therefore in digests. Even the shape of refusal is deterministic.
\end{remark}

\subsection{What a pass means}\label{ch09:pass-meaning}

The decision record's anatomy repays a moment's attention before its constraints. A decision carries the structural verdict; the constant-false authorization field; the code tuple; four optional digest fields identifying, where present, the action, the validation receipt, the compilation assessment, and the verified registry bytes that participated in the decision; and the Ring-4 boundary. The digest fields make every decision a provenance record: a reviewer holding a decision can name exactly which artifacts it adjudicated, fetch them by digest if retained, and re-run the gate to check that the same inputs still yield the same decision. Notably, the registry digest is populated \emph{only} when the registry's raw bindings verified---a decision never cites an unverified registry, so the citation itself carries information. Subclassing the decision type is rejected, and each digest field is syntax-validated, so a decision cannot be extended, impersonated, or made to cite a malformed identity.

The decision object that carries a pass is itself constrained so that it cannot be dressed up as more than it is. Its code tuple must be nonempty; a passing decision must carry an explicit structural-pass code and may carry no denial codes; a failing decision must carry an explicit denial code; and---a small rule with long reach---no code beginning with \texttt{ALLOW} is representable at all: the constructor rejects any decision containing one. The vocabulary of outcomes is \emph{pass/deny at the structural level} plus \emph{not granted at the authority level}, and the forbidden prefix ensures that no future refactor, log formatter, or downstream consumer can find an ``allowed'' in the kernel's output to seize upon, because the kernel cannot emit the word. Every decision, passing or failing, additionally carries the standing code \texttt{AUTHORIZATION\_NOT\_GRANTED\_RING4\_RESEARCH} and the Ring-4 boundary, so a decision record quoted out of context still states, in itself, what it is not.

\section{The two constant-false invariants}\label{ch09:constant-false}

\subsection{The invariants stated}\label{ch09:invariants-stated}

\begin{invariant}[No representable authorization]\label{ch09:inv:no-yes}
In v0.1: (i) the decision field \texttt{authorization\_granted} is constant false---the decision constructor raises if handed any other value, so a decision granting authorization is not a constructible object; and (ii) the backward-compatible spelling \texttt{allowed} is a computed property that returns false unconditionally---it reads no state, consults no policy, and has no writable backing field. There is no code path, in the gate or outside it, that yields a decision in which either is true.
\end{invariant}

Two corollaries give the invariant its teeth. First: \emph{claim-bearing intent is denied even against a locally constructed \textsc{licensed} ledger}. A caller may build a syntactically perfect claim registry whose records carry the text \textsc{licensed}, load it through the byte-backed loader, present it to the gate with a claim-bearing action---and the action denies with \texttt{DENY\_CLAIM\_BEARING\_ACTION\_NOT\_AUTHORIZED} exactly as it would with no registry at all. The registry chapter's slogan---\textsc{licensed} text inside a record is necessary metadata, not sufficient authority---is here made operational: registry content participates in integrity checks and identification, and participates in permission not at all, because there is no permission logic for it to participate in. Second: \emph{a self-hash is not issuance provenance}. Every receipt and descriptor in the kernel carries a canonical self-digest, and every such digest is scoped, in an explicit validated field, as corruption detection only. A verified self-digest proves the artifact is unmodified since construction; it says nothing about who constructed it, under what right, or with what review---and the assessment denies any receipt whose scope field claims otherwise (\texttt{DENY\_MISSTATED\_RECEIPT\_INTEGRITY\_SCOPE}). The historical incident behind this rule bears restating: the program's past contains AI-authored, self-hashed ``licensing receipts'' that circulated as if hash-validity were authority. In this kernel that artifact class is not merely rejected; it is \emph{addressed by name}, in a validated field, in every receipt.

\begin{example}[The licensed-ledger attempt, step by step]\label{ch09:ex:ledger}
It is instructive to trace the strongest locally mountable attempt at authorization through the machinery, because every step is one the kernel's history says someone will eventually take in good faith.

An operator---or an agent acting for one---prepares immaculate materials. A legacy source compiles cleanly through the full closure protocol; the compilation, receipt, and assessment all verify; a fresh replay would pass. A claim registry is constructed whose bytes are schema-valid, whose loader accepts them, and whose relevant record carries ring and license fields reading as licensed; the registry's raw bindings verify. An action descriptor is built with operation \texttt{PUBLISH\_CLAIM}, target \texttt{PUBLIC\_CLAIM\_SURFACE}, effect class \texttt{CLAIM\_BEARING}, and correctly paired claim identifier and version. Every digest in sight is genuine. Nothing is forged; every object is exactly what it says it is.

The gate's treatment: the assessment machinery, replaying the compilation, computes its findings against the action and reaches the effect-class arm, which yields \texttt{DENY\_CLAIM\_BEARING\_ACTION\_NOT\_AUTHORIZED}; the operation and target arms add their own denials, since only the legacy evaluation of a local Ring-4 artifact is admitted. The registry, verified, contributes its raw digest to the decision record---as context, nothing more. The decision emerges with \texttt{structurally\_admissible} false, \texttt{authorization\_granted} false (as it must: no other value is constructible), the denial codes, and the standing not-granted code. The registry's licensed text was read by nothing, weighed by nothing, and could have said anything.

Now weaken the attempt to its subtler form: the same materials, but the action claims effect class \texttt{PURE\_LOCAL\_COMPUTE} while carrying the claim metadata---a laundering shape, local-compute clothing over claim-bearing intent. This version dies earlier: the descriptor constructor rejects the combination outright, and the gate never sees an object. The two outcomes bracket the design: incoherent intent is unrepresentable; coherent claim-bearing intent is representable and constantly denied. Between them there is no seam.
\end{example}

\subsection{Why ``no code path to yes'' beats ``policy says no''}\label{ch09:no-path}

Both designs produce the same observable behavior today: every request is denied. The difference is in what must happen for tomorrow to differ, and in who can make it happen. The argument deserves to be made carefully, because it is the kernel's central design stance.

A policy-based denial is a \emph{data} property. Somewhere a configuration, a rule set, a ledger, or a flag encodes ``deny''; an evaluation engine reads it and complies. The security of the arrangement is the security of the data path: whoever can write the policy store, inject a higher-precedence rule, substitute the file the engine reads, exploit a parsing divergence between what auditors read and what the engine reads, or induce the engine to read the wrong policy at the wrong time, can flip the answer---without changing a line of reviewed code. The attack surface is every writer of the data, every reader ambiguity, and every cache and race in between. And policy has a second, quieter failure mode with no attacker at all: \emph{drift}. Policies are written to be changeable; changeable things change. An emergency exception added under deadline, a migration that resets a default, a well-meaning operator following a runbook---each is a plausible, local, reviewable-looking event, and their composition is a yes that no one ever decided.

A constant-false invariant is a \emph{structural} property. The gate does not evaluate ``should this be authorized?'' and conclude no; the gate has no authorization evaluation, and its output type has no inhabitant expressing yes. To make tomorrow differ, someone must change the code: alter the constructor that rejects true, add the evaluation logic that does not exist, introduce the trust-root machinery of Section~\ref{ch09:future}. Those changes are visible in exactly the places the program's review discipline watches---source diffs, test suites, the frozen contracts---and they cannot be made by any manipulation of inputs, configuration, environment, or stored state, because no input, configuration, environment, or state is consulted. In the language of security engineering, the design substitutes \emph{incapability} for \emph{restraint}: the classic distinction between a reference monitor that checks requests against policy and a system in which the dangerous operation is simply not implemented \cite{anderson}. Restraint requires the monitor to be correct on every path under every input forever; incapability requires the implementation to remain absent, which is a property a reviewer can confirm by reading and a test can confirm by exhaustion.

The testing asymmetry is worth spelling out, because it converts the philosophical point into engineering practice. The proposition ``the policy denies everything'' is a universally quantified claim over the policy-input product space---all requests times all reachable policy states---and testing samples it. The proposition ``\texttt{authorization\_granted} is constant false'' factors through the type: one test constructs a decision with the field true and asserts the constructor raises; one test greps the decision logic for any assignment of true; property-based tests then hammer the gate with arbitrary well-formed and malformed evidence and assert the field on every output. The Gate-4 suite does exactly this hammering, and the claim it evidences is the strong one: not ``we found no input that authorized'' but ``the output type rejected authorization in every constructed instance.'' A reviewer auditing the kernel for its most important property has a finite job.

There is also a compositional consequence. Systems embed kernels in wrappers, and wrappers have bugs. A wrapper around a policy-denying gate can misread a rich decision object---mistake structural admissibility for permission, surface an internal code as an outcome, cache a stale answer. A wrapper around this gate can certainly still misbehave, but the specific misbehavior of \emph{truthfully reporting an authorization} is unavailable, because no decision object carrying one exists to be misreported. The forbidden-\texttt{ALLOW} code rule and the standing not-granted code are the same idea applied to the decision's human-readable surface: the kernel's outputs are engineered so that even quotation out of context cannot manufacture a yes. What the constant-false design costs is flexibility---and Section~\ref{ch09:future} is the accounting of that cost: the flexibility is not abandoned but exiled, behind a list of preconditions that name exactly what a real yes would require.

\begin{threat}[The persuadable gate]\label{ch09:threat:persuasion}
The modern version of the policy-flipping attack deserves its own statement, because the CAS program operates AI agents as first-class participants in its own development. An agent tasked with ``getting the pipeline green'' or ``completing the release'' experiences a denial as an obstacle, and agents are competent obstacle-removers: they edit configuration, relax rules, construct satisfying inputs, and---the documented historical case---generate artifacts that look like the missing credential. Against a policy-based gate, every one of those moves is potentially sufficient, because the gate's answer supervenes on data the agent can write. Against this gate, the entire class is insufficient by construction: no configuration, ledger, receipt, environment variable, or crafted input is consulted for authorization, so no data-plane action changes the answer. The only remaining move is editing the kernel's code---which is precisely the move that the program's review discipline is built to catch, and which no longer resembles ``configuring the system'' even to a motivated reader. The constant-false design can be understood as moving the yes across the data/code boundary: from the plane agents and operators routinely touch to the plane where changes are diffs with reviewers. \emph{Blocked by} Invariant~\ref{ch09:inv:no-yes}, the forbidden-\texttt{ALLOW} rule, and the exact-type discipline of Rule~\ref{ch09:rule:typed-actions}.
\end{threat}

\begin{proposition}[Strength ordering]\label{ch09:prop:strength}
Let $D$ be the property ``every authorization request is denied under the deployed policy'' and $S$ the property ``no reachable code path constructs an authorization grant.'' Then $S$ implies $D$ for every policy state, every input, and every data-channel manipulation, while $D$ implies $S$ under no additional assumptions; and $S$, unlike $D$, is invariant under changes to configuration, stored state, and input distribution.
\end{proposition}

\begin{proof}[Design argument]
$S$ rules out grant construction structurally, so no policy state or input can yield a grant: $D$ holds under $S$ unconditionally. Conversely a system may satisfy $D$ today through policy content alone while retaining code paths that construct grants for policy states not currently deployed; such a system fails $S$ while satisfying $D$, so $D$ does not imply $S$. The invariance claim holds because $S$ quantifies over code, not over data: the witnesses to $S$---constructor rejection, absence of grant logic---are unaffected by any change that does not modify code.
\end{proof}

\section{The future-authorization design constraints}\label{ch09:future}

The kernel's constant no is not a claim that no should be the permanent answer. It is a claim about \emph{sequencing}: authorization is a capability that must be designed, reviewed, and admitted as a whole, and until then even a perfect structural record changes nothing. The v0.1 contract therefore does something unusual for a system with no authorization feature: it specifies, normatively, what any future implementation must bind before its first yes. Six bindings are named. This section develops each---what it is, why it is necessary, and what its absence would permit. The last question is the important one: each binding is best understood as the negation of a concrete laundering channel, several of which the program has already met in the wild \cite{shostack}.

\emph{First: an exact trusted registry and schema manifest.} The authorization decision must name, by digest, the exact registry bytes and schema bytes under which the claim it authorizes was recorded---not ``the claim registry'' as a role, but \emph{these} bytes. Necessity: the registry is where claim text, ring, and license state live; an authorizer that resolves the registry by name or path delegates its premise to whoever controls the resolution. Absence permits the founding attack of this kernel's threat history: a locally constructed ledger---syntactically valid, honestly loaded, saying \textsc{licensed}---becomes authoritative because nothing distinguishes it from the real one. v0.1 already refuses this even without an authorizer; a future authorizer without manifest binding would reintroduce it at the worst possible layer.

\emph{Second: an authenticated authority and trust root.} Every artifact the authorization consumes---the manifest above, the receipts below---must be verifiable against a cryptographic root of trust whose key custody lies outside the code that requests and outside the code that decides. Necessity: this is the one binding that no amount of structure can substitute for, and the reason v0.1 has no authorization at all---the kernel possesses hashing but not authentication, and Chapter~\ref{ch07:worker-protocol}'s walkthrough ended precisely at the question ``who vouches?''. Absence permits self-issuance in every form: any party able to run the kernel can mint every artifact the authorizer checks, and the self-hashed licensing receipts of the program's past return with better formatting. An unauthenticated authorization system is the constant-false invariant's negation with extra steps.

\emph{Third: reviewer and admission receipts.} The decision must bind signed records of the human review acts---who reviewed the claim, under what admission procedure, with what outcome---not merely the fact that a review state field says \textsc{approved}. Necessity: the program's claim discipline is anchored in accountable human judgment; the receipts are what make the judgment an artifact rather than an assertion, giving the authorization an audit trail that survives personnel and time. Absence permits state-flipping: whoever can write the review-state field \emph{becomes} the review, and a green test, a scripted migration, or an AI agent updating a ledger performs the promotion the kernel's entire history warns against. With receipts bound, a forged approval requires forging a signed artifact under the trust root; without them, it requires a one-line write.

\emph{Fourth: Evidence Packet identities and verification status.} The decision must bind the exact identities of the evidence packets supporting the claim \emph{and} their verification status---evidence checked, by whom, with what result---so the authorized claim is inseparable from its verified support \cite{prov}. Necessity: a claim is a conclusion over evidence; authorizing the conclusion while leaving the evidence unbound authorizes a sentence, not a finding. Absence permits evidence rot and evidence swapping: a claim authorized against strong packets continues to circulate after a packet is withdrawn, corrected, or shown unreproducible, or is quietly re-pointed at weaker evidence---and no consumer of the authorization can tell, because the authorization never said which evidence it meant. Binding identity plus verification status makes the authorization falsifiable in the only useful sense: invalidate the evidence and the authorization visibly no longer covers what it claims to cover.

\emph{Fifth: the claim identifier and exact version.} The decision must name one claim at one version, under the identifier discipline the descriptor already enforces. Necessity: claims evolve, and the difference between version $n$ and $n{+}1$ of a claim is frequently the difference between what review approved and what review would have rejected. Absence permits scope creep by versioning: an authorization issued for a modest v1 silently blankets an ambitious v3 because both share an identifier; or, with sloppy identifier discipline, blankets a different claim whose identifier normalizes to the same string. The registry chapter's rejection of ambiguous versions, and the descriptor's paired-fields rule, are this binding's existing foundations; the future authorizer inherits them or reopens the hole.

\emph{Sixth: the code, model, environment, semantic, and calibration identities of the requested operation.} The decision must bind, by digest, the operation that is actually about to run: which code, which model (or the canonical no-model identity), which execution environment, which semantic profile, and which calibration. Necessity: an authorization is a statement that a \emph{reviewed} operation may proceed; if the operation's identity is unbound, the statement attaches to a name, and names are re-pointable. Absence permits the time-of-review to time-of-use substitution that Chapter~\ref{ch07:worker-protocol} closed at the compiler level, reopened at the authority level: the reviewed model swapped for a fine-tuned variant, the reviewed environment for one with different dependencies, the reviewed calibration for none---all under the authorized name, all ``authorized.'' The compatibility profile's existing machinery---source digests, execution profiles, replay---exists so that this binding has something exact to bind; the authorizer completes the chain or the chain's strongest links dangle from nothing.

\begin{ruleenv}[Conjunction, and the default meanwhile]\label{ch09:rule:conjunction}
The six bindings are jointly necessary: an implementation lacking any one reduces, through that one's absence, to a laundering channel already catalogued. They are bound together in a separately reviewed design---reviewed as an authorization system, not assembled incrementally inside the research kernel---and until such a design is admitted in full, the kernel's answer to every claim-bearing request remains the constant of Invariant~\ref{ch09:inv:no-yes}: no.
\end{ruleenv}

It is worth noticing the relationship between the six bindings and the kernel that refuses to need them yet. Each binding demands an \emph{exact identity} for something---registry bytes, review acts, evidence packets, claim versions, operation components---and the kernel's five mechanisms are, in retrospect, the factories for exactly those identities: canonical encoding makes digests stable, registries make claim bytes exact, the compatibility protocol makes operation identity replayable, the boundary object keeps every artifact labeled while it waits. The constant-false gate is not a placeholder bolted onto an unfinished system; it is the correct decision function for a system that can already \emph{identify} everything and cannot yet \emph{authenticate} anything. The day the trust root exists, the yes it enables will have exact objects to bind---because the kernel spent v0.1 refusing to say yes to anything less.

\section{The denial vocabulary, collected}\label{ch09:denials}

As with the compatibility failure codes of Chapter~\ref{ch08:concurrency}, the decision model's denial vocabulary is a closed, typed set, and collecting it in one place (Table~\ref{ch09:tab:denials}) makes the model's coverage inspectable: every axis of the action vocabulary, every layer of the evidence hierarchy, and every binding of the receipt has a named refusal. The table also records the three non-denial codes---the single structural pass and its two mandatory limitations---so the complete output vocabulary of the gate is visible at once. It is a short list for a decision system, and the shortness is the point: a reviewer can hold the entire space of things this gate can say in one page, and verify that ``yes, with authority'' is not among them.

\begin{longtable}{@{}ll@{}}
\caption{Decision-model codes at v0.1.}\label{ch09:tab:denials}\\
\toprule
Code & Refuses / records \\
\midrule
\endfirsthead
\toprule
Code & Refuses / records \\
\midrule
\endhead
\bottomrule
\endfoot
\multicolumn{2}{@{}l}{\emph{Evidence-tier denials}} \\
\texttt{DENY\_MISSING\_COMPILATION\_ASSESSMENT} & no structural evidence supplied \\
\texttt{DENY\_UNKNOWN\_STRUCTURAL\_EVIDENCE} & evidence of unrecognized exact type \\
\texttt{DENY\_STRUCTURAL\_RECEIPT\_NOT\_AUTHORIZATION} & bare receipt presented as evidence \\
\texttt{DENY\_ASSESSMENT\_TYPE\_MISMATCH} & assessment component fails exact typing \\
\texttt{DENY\_CORRUPT\_COMPILATION\_ASSESSMENT} & assessment self-digest fails \\
\texttt{DENY\_ACTION\_ASSESSMENT\_MISMATCH} & supplied action differs from assessed action \\
\texttt{DENY\_COMPILATION\_REPLAY\_FAILED} & fresh replay could not be completed \\
\texttt{DENY\_RECEIPT\_REPLAY\_MISMATCH} & receipt differs from replay-derived receipt \\
\texttt{DENY\_COMPILATION\_EVIDENCE\_MISMATCH} & evidence digest differs from fresh digest \\
\texttt{DENY\_ASSESSMENT\_STATE\_MISMATCH} & stated verdict or codes differ from recomputation \\
\midrule
\multicolumn{2}{@{}l}{\emph{Receipt-finding denials}} \\
\texttt{DENY\_CORRUPT\_COMPILER\_VALIDATION} & receipt self-digest fails \\
\texttt{DENY\_COMPILER\_RESULT\_INVALID} & compiler reported an invalid result \\
\texttt{DENY\_UNKNOWN\_COMPILER\_RAW\_IDENTITY} & raw digest outside frozen variants \\
\texttt{DENY\_UNKNOWN\_COMPILER\_CONTENT\_IDENTITY} & LF digest not the frozen value \\
\texttt{DENY\_UNKNOWN\_COMPILER\_SOURCE\_AUTHORITY} & source authority not canonical \\
\texttt{DENY\_UNKNOWN\_SEMANTIC\_PROFILE} & semantic profile not the frozen string \\
\texttt{DENY\_UNKNOWN\_SYNTAX\_PROFILE} & syntax profile not the frozen string \\
\texttt{DENY\_NON\_PYTHON\_ENGINE\_PROFILE} & engine other than Python-only \\
\texttt{DENY\_NONISOLATED\_PYTHON\_PROCESS} & isolation booleans not satisfied \\
\texttt{DENY\_FALSE\_OS\_ISOLATION\_CLAIM} & receipt claims sandbox or host closure \\
\texttt{DENY\_NONLEGACY\_RECEIPT} & receipt not marked \texttt{legacy\_only} \\
\texttt{DENY\_MISSTATED\_RECEIPT\_INTEGRITY\_SCOPE} & self-hash scope overstated \\
\midrule
\multicolumn{2}{@{}l}{\emph{Action-binding denials}} \\
\texttt{DENY\_CORRUPT\_ACTION\_DESCRIPTOR} & descriptor self-digest fails \\
\texttt{DENY\_ACTION\_SOURCE\_MISMATCH} & action and receipt source digests differ \\
\texttt{DENY\_ACTION\_RESULT\_MISMATCH} & action and receipt result digests differ \\
\texttt{DENY\_ACTION\_ENVIRONMENT\_MISMATCH} & action and receipt environments differ \\
\texttt{DENY\_UNBOUND\_OR\_UNEXPECTED\_MODEL\_IDENTITY} & model identity not the no-model constant \\
\texttt{DENY\_UNSUPPORTED\_RUNTIME\_OPERATION} & operation not legacy evaluation \\
\texttt{DENY\_UNSUPPORTED\_RUNTIME\_TARGET} & target not local Ring-4 artifact \\
\texttt{DENY\_CLAIM\_BEARING\_ACTION\_NOT\_AUTHORIZED} & claim-bearing intent \\
\texttt{DENY\_EXTERNAL\_SIDE\_EFFECT\_NOT\_AUTHORIZED} & external side-effect intent \\
\texttt{DENY\_UNKNOWN\_EFFECT\_CLASS} & effect class outside the enum \\
\texttt{DENY\_UNKNOWN\_ACTION\_DESCRIPTOR} & action argument fails exact typing \\
\midrule
\multicolumn{2}{@{}l}{\emph{Registry denials}} \\
\texttt{DENY\_UNKNOWN\_CLAIM\_REGISTRY} & registry argument fails exact typing \\
\texttt{DENY\_CLAIM\_REGISTRY\_RAW\_BINDING} & registry raw byte bindings fail \\
\midrule
\multicolumn{2}{@{}l}{\emph{Non-denial codes}} \\
\texttt{STRUCTURAL\_PASS\_RING4\_PURE\_LOCAL\_COMPUTE} & the single structural pass \\
\texttt{LIMITATION\_NO\_OS\_SANDBOX} & mandatory limitation on every pass \\
\texttt{LIMITATION\_HOST\_OS\_DEPENDENCIES\_NOT\_CONTENT\_ADDRESSED} & mandatory limitation on every pass \\
\texttt{AUTHORIZATION\_NOT\_GRANTED\_RING4\_RESEARCH} & standing code on every decision \\
\end{longtable}

\section{Coda: the keystone property}\label{ch09:coda}

The decision model completes an arch the earlier chapters built. Semantic identity ensures the kernel cannot be confused about \emph{what} is named; the boundary object, about \emph{where} in the research lifecycle it stands; the encoding, about \emph{whether two things are the same}; the compatibility protocol, about \emph{what actually ran}; the registries, about \emph{what was actually recorded}. The decision model adds the keystone: the kernel cannot be confused---or talked, configured, or drifted---into saying \emph{yes}. Structural consistency is never authorization; the strongest artifact the kernel can produce is a fully replayed, fully validated, Ring-4, pure-local, not-granted structural pass that names its own limitations in typed fields.

That is a modest artifact by the standards of systems that promise more. The argument of this monograph has been that the modesty is the contribution: every stronger statement the program will ever responsibly make must be \emph{added}---through the six bindings, under review, over exactly identified objects---and the kernel's design ensures that until each addition is made, its absence is a typed refusal rather than a quiet assumption. The remaining chapters turn to the threat model that stress-tests this stance and to the evidence program that keeps its claims bounded.
